\BourbakiExercisesSection

¶ 1. 设 A 为交换环，E、F 为两个 A-代数，M 为左 E-模，N 为左 F-模，G 为代数 \(\mathbf{G} = \mathbf{E} \otimes_{\mathbf{A}} \mathbf{F}\) 。对每一对有序的自同态 \(u \in \operatorname{End}_{\mathbf{E}}(\mathbf{M})\) , \(v \in \operatorname{End}_{\mathbf{F}}(\mathbf{N})\) ，存在唯一的 G-自同态 \(w\) （定义在 \(\mathbf{M} \otimes_{\mathbf{A}} \mathbf{N}\) 上），使得 \(w(x \otimes y) = u(x) \otimes v(y)\) 对所有 \(x \in \mathbf{M}\) , \(y \in \mathbf{N}\) 成立，并且我们可写 \(w = \phi(u \otimes v)\)，其中 \(\phi\) 是从 \(\operatorname{End}_{\mathbf{E}}(\mathbf{M}) \otimes_{\mathbf{A}} \operatorname{End}_{\mathbf{F}}(\mathbf{N})\) 到 \(\operatorname{End}_{\mathbf{G}}(\mathbf{M} \otimes_{\mathbf{A}} \mathbf{N})\) 中的 A-线性映射（称为典范映射）。以下设 A 为一个域。

\begin{enumerate}
\def\labelenumi{(\alph{enumi})}
\item
  证明典范映射 \(\phi\) 是单射（考虑元素 \(z = \sum_{i} u_{i} \otimes v_{i}\) ，其中 \(v_{i} \in \operatorname{End}_{\mathbf{F}}(\mathbf{N})\) 在 A 上线性无关，并写出 \((\phi(z))(x_{\alpha} \otimes y) = 0\) ，其中 \(x_{\alpha}\) 取遍 M 在 A 上的一个基的元素，而 \(y \in \mathbf{N}\) 任意）。
\item
  假设 \(\mathbf{M}\) 是有限生成的自由 E-模。证明在下列每种情形映射 \(\phi\) 都是双射：（1）E 在域 A 上是有限秩的；（2）N 是有限生成的 F-模。
\item
  假设 \(\mathbf{M}\) 是自由 E-模，且存在 \(y \in \mathbf{N}\) 与自同态的无限序列 \((v_n)\) （诸自同态定义在 F-模 \(\mathbf{N}\) 上），使得 \(\mathbf{N}\) 的由诸 \(v_n(y)\) 生成的（A 上的）向量子空间在 A 上具有无限秩。证明：若映射 \(\phi\) 是双射，则 \(\mathbf{M}\) 在 \(\mathbf{E}\) 上的每个基都必然是有限的。
\item
  仍设 \(\mathbf{M}\) 是自由 E-模。证明：G-模 \(\mathbf{M} \otimes_{\mathbf{A}} \mathbf{N}\) 是忠实的，当且仅当 F-模 \(\mathbf{N}\) 是忠实的（考虑 \(\mathbf{E}\) 在 \(\mathbf{A}\) 上的一个基）.
\end{enumerate}

\begin{enumerate}
\def\labelenumi{\arabic{enumi}.}
\setcounter{enumi}{1}
\tightlist
\item
  设 K 是交换域，E、F 是两个中心包含
\end{enumerate}

K 的域，M 是 E 上的左向量空间，N 是 F 上的左向量空间；设 \(G = E \otimes_{K} F\) 。

\begin{enumerate}
\def\labelenumi{(\alph{enumi})}
\item
  设 \(x \neq 0\) 是 \(\mathbf{M}\) 的一个元素，\(y \neq 0\) 是 \(\mathbf{N}\) 的一个元素。证明 \(x \otimes y\) 在 G-模 \(\mathbf{M} \otimes_{\mathbb{K}} \mathbf{N}\) 中是自由的（利用习题 1(d) 以及 \(\operatorname{End}_{\mathbf{E}}(\mathbf{M})\)（相应地 \(\operatorname{End}_{\mathbf{F}}(\mathbf{N})\)）在 \(\mathbf{M}\)（相应地 \(\mathbf{N}\)）中异于 0 的元素所成集合上的传递性）.
\item
  证明 \(\mathbf{M} \otimes_{\mathbb{K}} \mathbf{N}\) 是自由 G-模。（若 \((m_{\alpha})\) 是 \(\mathbf{M}\) 在 \(\mathbf{E}\) 上的一个基，\((n_{\beta})\) 是 \(\mathbf{N}\) 在 \(\mathbf{F}\) 上的一个基，则用与 (a) 类似的方法证明诸 G-模 \(\mathrm{G}(m_{\alpha} \otimes n_{\beta})\) 之和是直和。）
\end{enumerate}

